% 6 Evaluation (~2.4 p).
\section{Evaluation}
\label{sec:evaluation}

\subsection{Research Questions}
\label{sec:rqs}
\begin{description}
  \item[RQ1 (effectiveness)] Under the same rewrite budget, does guided
    selection find more distinct inconsistencies than random and
    syntax-based selection?
  \item[RQ2 (cost)] What does each stage of the pipeline cost, and what does
    guidance add?
  \item[RQ3 (validity)] Can the variants be trusted to be equivalent to their
    originals?
\end{description}

\subsection{Subjects and Setup}
\label{sec:setup}
The subjects are \NumMethods{} JPAMB methods: \NumMethodsOrig{} from the
benchmark and \NumGuardsMethods{} methods (\NumGuardsCases{} annotated cases)
in a new class, \texttt{Guards}, written before any RQ1 measurement in five
patterns fixed in advance: guarded division, range checks before array
accesses, branches an analysis can decide, assertions on facts from an earlier
branch, and else-if chains used after the merge. Since the author of the
technique wrote them, every RQ1 result is also reported without them.

There are \NumSites{} eligible sites in \NumSiteMethods{} methods
(\NumSitesNoGuards{} without \texttt{Guards}). \NumUncertain{}
(\UncertainPercent\%) are candidates for \texttt{guided}, \NumStrict{} of them
also strict candidates for \texttt{guided-strict}; \NumSettled{} are settled and \NumNoState{} have no
guide state. The share of candidates is high because many conditions compare
unconstrained parameters, array accesses are always uncertain, and sites on one
line share a flag. The budget is $N=\BudgetHalf$ (50\% of the eligible sites)
or $N=\BudgetQuarter$ (25\%), and $N=\BudgetHalfNoGuards$ or
$N=\BudgetQuarterNoGuards$ without \texttt{Guards}, with \NumSeeds{} seeds per
strategy. The measure is the number of distinct inconsistencies per run,
summed over the three analyzers; the tests use the sum because
\texttt{static-interval} has none, and \cref{tab:rq1} gives the per-analyzer
means. Strategies are compared by a two-sided Mann--Whitney U test (normal
approximation with tie and continuity correction) and the Vargha--Delaney effect size
$\hat{A}_{12}$, the probability that a \emph{run} (one strategy, budget and
seed) of the first strategy finds more
than a run of the second, ties counted half. All runs used one machine:
\Machine.

\subsection{RQ1: Effectiveness}
\label{sec:rq1}
Rewriting \emph{every} eligible site bounds what any strategy can find; we
call these inconsistencies the \emph{pool}. It yields only \Pool{} distinct inconsistencies (\PoolSign{} for
\texttt{static-sign}, \PoolSignNarrow{} for \texttt{static-sign-narrow},
\PoolInterval{} for \texttt{static-interval}), from \NumInconsistentVariants{}
of the \NumValid{} variants, all from R2-strict; without \texttt{Guards}, the
pool is \PoolNoGuards{}. RQ1 therefore asks whether a strategy reaches these
few sites within its budget.

\begin{table}[t]
\caption{RQ1: distinct inconsistencies per run at the 50\% budget, mean
(population standard deviation) over \NumSeeds{} seeds, per analyzer and in total, with
the Mann--Whitney $p$ and $\hat{A}_{12}$ of the total against
\texttt{random}. $^\dagger$Exploratory, defined after the pool was known.}
\label{tab:rq1}
\small
% Generated by experiments/paper_numbers.py -- do not edit by hand.
\begin{tabular}{@{}lrrrrrr@{}}
\toprule
Strategy & \texttt{sign} & \texttt{sign-narrow} & \texttt{interval} & total & $p$ vs.\ random & $\hat{A}_{12}$ \\
\midrule
\multicolumn{7}{@{}l}{\emph{all: 408 sites, $N=204$, pool 5}} \\
\texttt{random} & 1.1 (0.8) & 2.0 (1.0) & 0.0 (0.0) & 3.1 (1.8) &  &  \\
\texttt{syntax} & 1.3 (0.5) & 2.3 (0.5) & 0.0 (0.0) & 3.6 (0.9) &  &  \\
\texttt{guided} & 0.0 (0.0) & 0.3 (0.5) & 0.0 (0.0) & 0.3 (0.5) & 0.002 & 0.10 \\
\texttt{guided-strict} & 0.0 (0.0) & 0.0 (0.0) & 0.0 (0.0) & 0.0 (0.0) & $<0.001$ & 0.05 \\
\texttt{settled}$^\dagger$ & 2.0 (0.0) & 3.0 (0.0) & 0.0 (0.0) & 5.0 (0.0) & 0.006 & 0.80 \\
\midrule
\multicolumn{7}{@{}l}{\emph{without Guards: 344 sites, $N=172$, pool 4}} \\
\texttt{random} & 1.1 (0.8) & 1.1 (0.8) & 0.0 (0.0) & 2.2 (1.7) &  &  \\
\texttt{syntax} & 1.3 (0.5) & 1.3 (0.5) & 0.0 (0.0) & 2.6 (0.9) &  &  \\
\texttt{guided} & 0.0 (0.0) & 0.0 (0.0) & 0.0 (0.0) & 0.0 (0.0) & 0.002 & 0.15 \\
\texttt{guided-strict} & 0.0 (0.0) & 0.0 (0.0) & 0.0 (0.0) & 0.0 (0.0) & 0.002 & 0.15 \\
\texttt{settled}$^\dagger$ & 2.0 (0.0) & 2.0 (0.0) & 0.0 (0.0) & 4.0 (0.0) & 0.006 & 0.80 \\
\bottomrule
\end{tabular}

\end{table}

\Cref{tab:rq1} gives the result. The main strategy,
\texttt{guided}, and its sensitivity variant \texttt{guided-strict} find
\emph{fewer} inconsistencies than \texttt{random} in every setting. At $N=\BudgetHalf$,
\texttt{guided} finds \MeanGuidedHalf{} on average, against \MeanRandomHalf{}
for \texttt{random} (\PGuidedRandomHalf, $\hat{A}_{12}=\AGuidedRandomHalf$) and
\MeanSyntaxHalf{} for \texttt{syntax} (\PGuidedSyntaxHalf); at
$N=\BudgetQuarter$ and without \texttt{Guards} the direction against
\texttt{random} is the same. Of
the \NumPlannedTests{} planned tests that can be computed, the largest $p$ is
\MaxPlannedP{}, below the Bonferroni threshold of \BonferroniAlpha{}; in the
\NumPlannedNoTest{} remaining comparisons, without \texttt{Guards} at
$N=\BudgetQuarterNoGuards$, \texttt{syntax} finds none either. These results
are largely fixed by construction. Because $N$ is smaller than the number of
candidates, \texttt{guided} never reaches a settled site, whatever the seed,
and \NumInconsistentSitesSettled{} of the \NumInconsistentVariants{} sites
that produce an inconsistency are settled; \PoolFromSettledOnly{} of the
\Pool{} distinct inconsistencies come from no uncertain site. The seeds vary
only which of these sites \texttt{random} and \texttt{syntax} reach.

\emph{Answer to RQ1: no. Choosing uncertain sites first finds fewer distinct
inconsistencies than random selection in all four settings, and fewer than
syntax-based selection in three; in the fourth, neither finds any.}

\paragraph{Exploratory: settled sites first}
The opposite order, which we defined only after seeing the pool and therefore
treat as a hypothesis, rewrites the \NumSettled{} settled sites first. Since
\NumSettled{} is below every budget, it contains the whole pool at both
budgets, by construction (\MeanSettledHalf{} in every seed). Whether settled
sites are richer elsewhere would need a new experiment on other analyzers and
programs.

\subsection{RQ2: Cost}
\label{sec:rq2}
One complete run of the pipeline, which reproduces the stored sites,
selections, validation results and outcome sets exactly,
takes \CostTotal{}\,s. Validation dominates with \CostValidation{}\,s:
compiling every variant and running every input on both versions. Computing
the guide's states for all methods takes \CostExport{}\,s, site finding and
selection \CostSites{}\,s and \CostSelection{}\,s, and analyzing every
original and variant with all three analyzers \CostAnalysis{}\,s. Guidance is
thus negligible, and under a time budget every strategy would validate about
the same number of variants, so a time-budget comparison reduces to the
rewrite-budget comparison of RQ1. For the same reason we do not report time to
the first inconsistency, which here is the position of the first of
\NumInconsistentVariants{} sites in each order.

\emph{Answer to RQ2: guidance costs \CostExport{}\,s of \CostTotal{}\,s;
validating the variants dominates the cost.}

\subsection{RQ3: Validity of Variants}
\label{sec:rq3}
All \NumCompile{} variants compile and agree with their originals on every
input, so testing discards \NumDiscarded{}, and \NumBytecodeIdentical{} have
the original's bytecode. To check that the inputs could have exposed a wrong rewrite, we
validated deliberately wrong ones: off-by-one or opposite operators for R2,
and negation without swapping the branches for R1. Of the \NumMutants{}
mutants, \NumMutantsEquivalent{} are equivalent to the original on all inputs
(for example, an off-by-one at a boundary the operand never reaches), each
with a written argument; testing rejects \NumMutantsKilled{} of the other
\NumMutantsNonEquivalent{}. The \NumMutantsSurvived{} survivors need
coordinated values in two or three parameters, which one-at-a-time synthesis
does not produce. For conditions on parameters the method never assigns, we
measured the original condition at entry: for \NumCondBoth{} variants the
inputs make it take both values and for \NumCondOne{} only one; for the other
\NumCondUnknown{} variants this is not measured.

\emph{Answer to RQ3: every variant passed, and the inputs reject all but
\NumMutantsSurvived{} non-equivalent wrong rewrites; the evidence is weaker for
methods with non-integer parameters (\cref{sec:threats}).}

\subsection{What the Inconsistencies Reveal}
\label{sec:findings}
\Cref{tab:inconsistencies} lists the variants that change an outcome set. They
share one cause: the sign domain abstracts the constant $0$ exactly, to
$\{0\}$, but $1$ to $\{+\}$ and $-1$ to $\{-\}$. A value in $\{+\}$ compared
with $\{0\}$ by \texttt{>} is always greater; compared with $\{+\}$ by
\texttt{>=}, it may not be, since the domain cannot tell $1$ from $2$. All
\NumInconsistentSitesCrossZero{} sites are among the \NumStrictCrossZero{}
R2-strict sites whose literal changes to or from $0$.

\begin{table}[t]
\caption{The variants with an inconsistency. Guide: the site's status for the
interval guide. Outcomes: ok, assertion error (AE), divide by zero (DZ),
non-termination ($\ast$).}
\label{tab:inconsistencies}
\small
% Generated by experiments/paper_numbers.py -- do not edit by hand.
\begin{tabular}{@{}lllll@{}}
\toprule
Site & Rewrite & Guide & Analyzers & Original $\to$ variant \\
\midrule
\texttt{Guards:138} & \texttt{x >= 0} $\to$ \texttt{x > -1} & uncertain & \texttt{sign-narrow} & \{ok\} $\to$ \{ok, AE\} \\
\texttt{Guards:139} & \texttt{y > 0} $\to$ \texttt{y >= 1} & settled & \texttt{sign-narrow} & \{ok\} $\to$ \{ok, AE\} \\
\texttt{Loops:19} & \texttt{i > 0} $\to$ \texttt{i >= 1} & settled & \texttt{sign}, \texttt{sign-narrow} & \{$\ast$\} $\to$ \{AE, $\ast$\} \\
\texttt{Loops:27} & \texttt{i > 0} $\to$ \texttt{i >= 1} & settled & \texttt{sign}, \texttt{sign-narrow} & \{$\ast$\} $\to$ \{DZ, $\ast$\} \\
\bottomrule
\end{tabular}

\end{table}

\begin{lstlisting}[float=t,caption={\texttt{Loops.neverAsserts} (left) and
its R2-strict variant (right).},label={lst:loops}]
int i = 1;                        int i = 1;
while (i > 0) {}                  while (i >= 1) {}
assert false;                     assert false;
\end{lstlisting}

In \cref{lst:loops}, both sign analyzers know $i\in\{+\}$. On the left the
loop condition is always true, the assertion is unreachable, and the
prediction is $\{\ast\}$; on the right the condition may be false and an
assertion error is added. \texttt{Loops.neverDivides} is the same with a
division by zero. The interval guide knows $i\in[1,1]$ and decides both
forms, so the site is settled and \texttt{guided} never reaches it within the
budget. \Cref{sec:overview} shows the same effect in a guard.

This explains why the hypothesis behind guidance (\cref{sec:sites}) fails for
these analyzers. The guide represents every constant exactly, the sign domain
only $0$. An inconsistency of this kind needs a comparison that the sign
analyzer decides in one form and not in the other, and such a comparison is
usually decided by the guide in both, so the site is settled. Where the guide is
uncertain, the condition is typically undecided for all three analyzers in
both forms, and the rewrite changes nothing; only
\NumInconsistentSitesUncertain{} of the \NumUncertain{} uncertain sites
produces an inconsistency.

The other rules change no outcome set on JPAMB. On each edge, both domains
refine the two operands of a comparison alike, whichever side they are on and
whichever edge is taken, so swapping operands or inverting a branch gives the
same edge states. For \texttt{static-sign-narrow}, splitting
\texttt{x != y} gives the union of the \texttt{<} and \texttt{>} refinements,
which is the \texttt{!=} refinement, and splitting \texttt{==} gave the same
states as \texttt{==} in every case we checked; for intervals the hull of the two parts equals the
\texttt{!=} refinement, and the sequence of \texttt{<=} and \texttt{>=} the
\texttt{==} refinement. For intervals the fixpoint could still depend on the
order in which a widening point is reached when R1 or R2-split reorders or
adds branches inside a loop, but this does not occur on JPAMB, and R2-strict
is exact for intervals, since $x>c$ and $x\ge c+1$ refine to the same
interval.
